\subsection{Experimental Setup}
\label{sec:experiment-setup}

The laser ultrasonic testing system comprised two primary components: an ultrasonic wave generation subsystem and a wave detection subsystem. For ultrasonic wave generation, a Q-switched Nd:YAG laser operating at its fundamental wavelength ($\lambda$ = 1064\,nm) was employed. The laser delivered pulses with duration on the order of nanoseconds, enabling the generation of broadband ultrasonic waves via the thermoelastic effect at the specimen surface. For wave detection, a laser Doppler vibrometer (LDV) was used to measure out-of-plane velocity at the specimen surface. In selected configurations, a piezoelectric transducer (PZT) with appropriate bandwidth was used as an alternative detection method.

The scanning configuration employed a computer-controlled translation stage to move either the generation laser or the specimen relative to the detection probe. Two scanning modes were evaluated: point-by-point mechanical scanning for discrete measurement locations, and line scanning along predefined paths across the specimen surface.

Two aluminum plate specimens were used in this study. The training specimen was an intact aluminum plate with nominal dimensions of 300\,mm $\times$ 300\,mm $\times$ 3\,mm (length $\times$ width $\times$ thickness). The testing specimen consisted of an aluminum plate of the same nominal dimensions containing a series of machined slots and notches of varying depths and widths, serving as controlled defect targets for evaluating detection capability.

Data acquisition was performed using a digital oscilloscope or equivalent high-speed acquisition system. The sampling rate was set to a minimum of twice the effective bandwidth to satisfy the Nyquist criterion. The acquisition bandwidth was configured to capture the full frequency content of the ultrasonic signals. Signal averaging was applied to improve the signal-to-noise ratio, with the number of averages adjusted based on the desired noise reduction and acquisition time constraints.

Specific parameters—including laser pulse energy, spot size at the specimen surface, scanning step size, and precise detection bandwidth—are provided as placeholders in the system configuration and will be finalized prior to experimental execution.

% Placeholder: laser pulse energy (e.g., 50--200\,mJ)
% Placeholder: laser spot diameter (e.g., 3--6\,mm)
% Placeholder: scanning step size (e.g., 1--5\,mm)
% Placeholder: LDV detection bandwidth (e.g., DC to 20\,MHz)
% Placeholder: number of averages (e.g., 32, 64, 128)